% Converted from the audited Markdown manuscript; responsible author: Carptopus.
\documentclass[11pt]{article}
\usepackage[a4paper,margin=28mm]{geometry}
\usepackage[T1]{fontenc}
\usepackage[utf8]{inputenc}
\usepackage{lmodern}
\usepackage{microtype}
\usepackage{amsmath,amssymb,amsthm,mathtools}
\usepackage{xcolor}
\usepackage[colorlinks=true,linkcolor=blue!55!black,citecolor=blue!55!black,urlcolor=blue!65!black]{hyperref}
\usepackage{fancyhdr}

\newtheorem{theorem}{Theorem}[section]
\newtheorem{lemma}[theorem]{Lemma}
\newtheorem{corollary}[theorem]{Corollary}
\numberwithin{equation}{section}

\pagestyle{fancy}
\fancyhf{}
\fancyhead[L]{Carptopus}
\fancyhead[R]{Complete periodic quaternion ringsets}
\fancyfoot[C]{\thepage}
\setlength{\headheight}{14pt}
\setlength{\parindent}{0pt}
\setlength{\parskip}{0.45em}
\setlength{\emergencystretch}{4em}

\title{Periodic ringsets in the Lipschitz quaternions:\\
the complete six-letter classification}
\author{Carptopus\\\href{mailto:carptopus@163.com}{carptopus@163.com}}
\date{24 August 2026}

\hypersetup{
  pdftitle={Periodic ringsets in the Lipschitz quaternions: the complete six-letter classification},
  pdfauthor={Carptopus},
  pdfsubject={A complete ringset classification for six-letter periodic subsets of the Lipschitz quaternions},
  pdfkeywords={integer-valued polynomial, Lipschitz quaternion, ringset, null ideal, periodic word, fixed divisor, split prime, ramified prime}
}

\begin{document}
\maketitle

\begin{abstract}

Let

\nopagebreak[4]
\[
\mathbf L=\mathbb Z[\mathbf i,\mathbf j,\mathbf k]
\]

be the ring of Lipschitz quaternions, and let

\nopagebreak[4]
\[
U=\{\pm\mathbf i,\pm\mathbf j,\pm\mathbf k\}
\]

be its six square roots of $-1$. For a periodic word
$w:\mathbb Z\to U$, we study

\nopagebreak[4]
\[
S_w=\{a+w(a):a\in\mathbb Z\}\subseteq\mathbf L.
\]

A subset of $\mathbf L$ is a ringset if the rational quaternion polynomials
integer-valued on it are closed under multiplication. We give a complete
criterion for $S_w$ to be a ringset for every prescribed period, including
even periods and words containing antipodal letters.

At ramified and nonsplit prime divisors of the period, every highest
prime-power position class must contain at least two letters. At a split
prime, the only obstruction occurs when two position fibres separated by a
square root of $-1$ are singletons carrying antipodal letters. Necessity is
proved by explicit null polynomials. The ramified prime $2$ requires a
Gaussian valuation calculation and a squared central fixed-divisor factor;
the split obstruction uses rank-one idempotents over a matrix ring. These
local results, periodic congruence orbits, and the Chinese remainder theorem
give the global if-and-only-if classification. The theorem strictly
generalizes the preceding squarefree-period and all-odd-period binary
classifications.

\end{abstract}

\section{Introduction}


Let $\mathbb D$ be the rational Hamilton quaternion algebra and let

\nopagebreak[4]
\[
\mathbf L
=\mathbb Z[\mathbf i,\mathbf j,\mathbf k],
\qquad
\mathbf i^2=\mathbf j^2=-1,
\qquad
\mathbf i\mathbf j=\mathbf k=-\mathbf j\mathbf i.
\]

The variable in $\mathbb D[x]$ is central. Polynomials have coefficients on
the left and are evaluated on the right:

\nopagebreak[4]
\[
\left(\sum_d c_dx^d\right)(\alpha)=\sum_dc_d\alpha^d.
\]

For $S\subseteq\mathbf L$, put

\nopagebreak[4]
\[
\mathrm{Int}(S,\mathbf L)
=\{f\in\mathbb D[x]:f(S)\subseteq\mathbf L\}.
\]

Following Werner \cite{werner2026}, $S$ is called a \emph{ringset} if
$\mathrm{Int}(S,\mathbf L)$ is a subring of $\mathbb D[x]$. This
condition is nontrivial because the coefficient algebra is noncommutative.

Werner \cite{werner2026} classified finite ringsets in the Lipschitz and Hurwitz
quaternion rings and supplied residue-null-ideal criteria for infinite sets.
The first paper in the present series \cite{carptopus-squarefree} classified periodic binary sets of
odd squarefree period. The second \cite{carptopus-odd} removed squarefreeness and classified
all odd-period words with alphabet $\{\mathbf i,\mathbf j\}$. Two natural
boundaries remained:

\begin{enumerate}
\item the ramified prime $2$, which is unavoidable for even periods;
\item the other four Lipschitz square roots of $-1$, particularly antipodal
pairs at split primes.

\end{enumerate}

This paper closes both boundaries. Let

\nopagebreak[4]
\[
U=\{\pm\mathbf i,\pm\mathbf j,\pm\mathbf k\},
\]

let $m\ge1$, and let $w:\mathbb Z\to U$ satisfy $w(a+m)=w(a)$. The displayed
period is not required to be least. Define

\nopagebreak[4]
\[
S_w=\{a+w(a):a\in\mathbb Z\}.
\]

For a prime $p\mid m$, write $\nu_p=v_p(m)$. For
$c\in\mathbb Z/p^{\nu_p}\mathbb Z$, define the position-fibre alphabet

\nopagebreak[4]
\[
A_p(c)
=\{w(a):a\equiv c\pmod {p^{\nu_p}}\}\subseteq U.
\]

Our main result is the following.

\begin{theorem}
\label{thm:main}

Let $w:\mathbb Z\to U$ have period $m\ge1$. Then $S_w$ is a ringset in
$\mathbf L$ if and only if all of the following hold:

\begin{enumerate}
\item $w$ is nonconstant;
\item for every prime $p\mid m$ such that either $p=2$ or
$p\equiv3\pmod4$, one has

$$
|A_p(c)|\ge2
\qquad
\left(c\in\mathbb Z/p^{\nu_p}\mathbb Z\right);
$$

\item for every prime $p\mid m$ with $p\equiv1\pmod4$, every
$s\in\mathbb Z/p^{\nu_p}\mathbb Z$ satisfying $s^2=-1$, every
$c\in\mathbb Z/p^{\nu_p}\mathbb Z$, and every $u\in U$, it is not the
case that

$$
A_p(c-s)=\{u\}
\quad\text{and}\quad
A_p(c+s)=\{-u\}.
$$

\end{enumerate}

\end{theorem}

Condition 3 is independent of the choice between the two square roots:
replacing $s$ by $-s$ interchanges the fibres and replaces $u$ by $-u$.
All three conditions are independent of the choice of displayed period.

For the binary alphabet $\{\mathbf i,\mathbf j\}$, condition 3 is automatic,
and Theorem 1.1 reduces on odd periods to \cite[Theorem 1.1]{carptopus-odd}. The new mechanisms
are the ramified $2$-adic obstruction and the exact antipodal obstruction at
split primes. The paired-fibre lemma, matrix representation, fixed-divisor
background, and residue-null-ideal bridge are not claimed as new.

The rest of the paper proves Theorem 1.1. Sections 2 and 3 collect the global
and paired-fibre tools. Section 4 isolates the split-prime geometry. Section
5 records a prime-power fixed-divisor estimate. Sections 6 and 7 construct
the nonsplit and ramified obstruction polynomials. Section 8 assembles the
classification.

\section{Residue null ideals and periodic fibres}


For a ring $A$ and $T\subseteq A$, define

\nopagebreak[4]
\[
\mathcal N(T,A)
=\{f\in A[x]:f(t)=0\text{ for every }t\in T\}.
\]

This is always a left ideal. The set $T$ is \emph{core} if
$\mathcal N(T,A)$ is a two-sided ideal. For $n\ge2$, let $\overline S_n$
denote the image of $S\subseteq\mathbf L$ in
$\mathbf L/n\mathbf L$.

We use the following consequence of Werner \cite[Lemmas 5.3 and 5.4]{werner2026}.

\begin{lemma}[Residue criterion]
\label{lem:residue}

Let $S\subseteq\mathbf L$.

\begin{enumerate}
\item If $\overline S_n$ is core in $\mathbf L/n\mathbf L$ for every $n\ge2$,
then $S$ is a ringset.
\item Conversely, if for some $n$ there are
$g\in\mathcal N(\overline S_n,\mathbf L/n\mathbf L)$ and
$r\in\mathbf L/n\mathbf L$ such that the coefficientwise right multiple
$gr$ is not null on $\overline S_n$, then $S$ is not a ringset.

\end{enumerate}

\begin{proof}

Write a rational quaternion polynomial as $f=g/n$ with
$g\in\mathbf L[x]$. By \cite[Lemma 5.3]{werner2026}, it is integer-valued on $S$ exactly
when the reduction of $g$ is null on $\overline S_n$. Stability of every
null ideal under coefficientwise right multiplication by constants gives
multiplicative closure by \cite[Lemma 5.4]{werner2026}. Conversely, a failed right multiple
lifts to an integer-valued polynomial whose corresponding product is not
integer-valued.
\end{proof}

\end{lemma}


The alphabet $U$ is intrinsic to the Lipschitz order.

\begin{lemma}
\label{lem:null-factor}

The elements of $\mathbf L$ whose square is $-1$ are exactly
$\pm\mathbf i,\pm\mathbf j,\pm\mathbf k$.

\begin{proof}

If $q=a+b\mathbf i+c\mathbf j+d\mathbf k$ and $q^2=-1$, its vector part is
$2a(b\mathbf i+c\mathbf j+d\mathbf k)$, so either $a=0$ or the vector part
vanishes. The latter is incompatible with $q^2=-1$ for integral $a$. Thus
$a=0$ and $b^2+c^2+d^2=1$, whose integral solutions are the six stated
units.
\end{proof}

\end{lemma}


We next record the congruence orbit of a periodic word.

\begin{lemma}[Periodic orbit lemma]
\label{lem:periodic-orbit}

Let $w$ have period $m$, let $q\ge1$, and fix
$a\in\mathbb Z/q\mathbb Z$. The positions modulo $m$ attained by integers
congruent to $a$ modulo $q$ are

\nopagebreak[4]
\[
a+q\mathbb Z\pmod m
=\{r\in\mathbb Z/m\mathbb Z:r\equiv a
  \pmod {\gcd(q,m)}\}.
\]

\begin{proof}

The subgroup generated by $q$ in $\mathbb Z/m\mathbb Z$ consists exactly
of the multiples of $\gcd(q,m)$. Translate by $a$.
\end{proof}

\end{lemma}


For $q=p^e$ and $\nu=v_p(m)$, a real residue fibre modulo $p^e$ therefore
sees one position class modulo $p^{\min(e,\nu)}$. In particular, a
highest-level class modulo $p^\nu$ controls every finer real residue fibre,
while a coarser fibre is a union of highest-level classes.

The same formula proves invariance under a nonminimal displayed period. If
$d\mid m$ is the least period, then a class modulo $p^{v_p(m)}$ visits
exactly a class modulo $p^{v_p(d)}$ when $p\mid d$, and visits the entire
least period when $p\nmid d$. The letter sets in Theorem 1.1 therefore
repeat, rather than change, when the displayed period is enlarged. At split
primes, Hensel lifting transports the two square-root offsets compatibly.

\section{More than one letter above every real residue}


Every point $a+u$ with $u\in U$ has central minimal polynomial

\nopagebreak[4]
\[
\mu_a(x)=(x-a)^2+1.
\]

The following paired-fibre lemma extends the calculation used in \cite{carptopus-squarefree,carptopus-odd}.

\begin{lemma}[Multi-letter paired fibres]
\label{lem:paired-fibres}

Let $n\ge2$, and let

\nopagebreak[4]
\[
T\subseteq\{a+u:a\in\mathbb Z/n\mathbb Z,\ u\in U\}.
\]

If at least two distinct letters occur above every real residue, then $T$ is
core in $\mathbf L/n\mathbf L$.

\begin{proof}

Take $f\in\mathcal N(T,\mathbf L/n\mathbf L)$ and fix a real residue $a$.
Divide $f$ by the central monic polynomial $\mu_a$. The remainder is

\nopagebreak[4]
\[
r(x)=\gamma_1x+\gamma_0.
\]

If $a+u$ and $a+v$ are two points of $T$ with $u\ne v$, subtraction of
their zero evaluations gives

\nopagebreak[4]
\[
\gamma_1(u-v)=0.
\]

If $v\ne-u$, then $\mathrm{Nm}(u-v)=2$, so right multiplication by
$\overline{u-v}$ yields $2\gamma_1=0$. If $v=-u$, then
$u-v=2u$, and right multiplication by the unit $-u$ gives the same
conclusion.

For a basis unit $t\in\{\mathbf i,\mathbf j,\mathbf k\}$ and
$\alpha=a+u$, direct right evaluation gives

\nopagebreak[4]
\[
(rt)(\alpha)-r(\alpha)t
=\gamma_1(t\alpha-\alpha t).
\]

Every displayed commutator lies in $2\mathbf L/n\mathbf L$, so the
right-hand side vanishes. Centrality of $\mu_a$ passes this conclusion from
the remainder to $f$. Thus the null ideal is stable under coefficientwise
right multiplication by the quaternion basis and hence by every constant.
Since it is a left ideal and $x$ is central, it is then stable under right
multiplication by every polynomial.
\end{proof}

\end{lemma}


Lemma 3.1 immediately supplies all ramified and nonsplit sufficiency
directions in Theorem 1.1. The remaining sufficiency issue is the split
case, where a single letter above a real residue can already suffice.

\section{Split primes and antipodal fibres}


Let $p\equiv1\pmod4$, let $e\ge1$, put

\nopagebreak[4]
\[
R=\mathbb Z/p^e\mathbb Z,
\qquad
A=\mathbf L/p^e\mathbf L,
\]

and choose $s\in R$ with $s^2=-1$. The standard representation

\nopagebreak[4]
\[
\mathbf i\longmapsto
\begin{pmatrix}s&0\\0&-s\end{pmatrix},
\qquad
\mathbf j\longmapsto
\begin{pmatrix}0&1\\-1&0\end{pmatrix}
\]

identifies $A$ with $M_2(R)$. For $u\in U$, define

\nopagebreak[4]
\[
e_u^+=\frac{1+u/s}{2},
\qquad
e_u^-=\frac{1-u/s}{2}.
\]

They are complementary rank-one idempotents, and

\nopagebreak[4]
\[
(a+u)^d=(a+s)^d e_u^+ +(a-s)^d e_u^-.
\]

Generators for their image lines may be chosen as follows.

\begin{center}
\renewcommand{\arraystretch}{1.08}
\begin{tabular}{c c c}
\hline
$u$ & $\mathrm{im}(e_u^+)$ & $\mathrm{im}(e_u^-)$ \\
\hline
$\mathbf i$ & $(1,0)^T$ & $(0,1)^T$ \\
$-\mathbf i$ & $(0,1)^T$ & $(1,0)^T$ \\
$\mathbf j$ & $(1,s)^T$ & $(1,-s)^T$ \\
$-\mathbf j$ & $(1,-s)^T$ & $(1,s)^T$ \\
$\mathbf k$ & $(1,1)^T$ & $(1,-1)^T$ \\
$-\mathbf k$ & $(1,-1)^T$ & $(1,1)^T$ \\
\hline
\end{tabular}
\end{center}

The determinants obtained from distinct same-sign lines are units. A plus
line for $u$ and a minus line for $v$ fail to be complementary exactly when
$v=-u$. Indeed, all nonzero determinants are among

\nopagebreak[4]
\[
\pm1,\quad \pm s,\quad \pm2,\quad \pm2s,
\quad \pm(1+s),\quad \pm(1-s),
\]

and these are units because $p$ is odd and
$(1+s)(1-s)=2$.

\begin{lemma}[Split local core criterion]
\label{lem:split-core}

Let

\nopagebreak[4]
\[
T\subseteq\{a+u:a\in R,\ u\in U\}
\]

have a nonempty fibre above every $a\in R$. Suppose that, for every
$z\in R$, the two fibres above $z-s$ and $z+s$ are not respectively the
singleton pair $\{u\}$ and $\{-u\}$ for any $u\in U$. Then $T$ is core in
$A$.

\begin{proof}

Take a null polynomial $f(x)=\sum_dc_dx^d$ and introduce the scalar-variable
matrix polynomial

\nopagebreak[4]
\[
F(z)=\sum_dc_dz^d.
\]

At fixed $z$, every letter $u$ occurring above $z-s$ gives

\nopagebreak[4]
\[
F(z)e_u^+=0,
\]

and every letter $v$ occurring above $z+s$ gives

\nopagebreak[4]
\[
F(z)e_v^-=0.
\]

If either fibre contains two letters, the corresponding two same-sign image
lines are complementary. If both are singletons, the hypothesis and the
table show that the plus and minus lines are complementary. Thus $F(z)=0$
for every $z\in R$.

For any $B\in A$, the scalar polynomial associated with the coefficientwise
right multiple $fB$ is $F(z)B=0$. The same idempotent decomposition shows
that $fB$ vanishes on $T$. Hence the null ideal is stable under all constant
right multiples and therefore is two-sided.
\end{proof}

\end{lemma}


The excluded singleton-antipode configuration is also a genuine periodic
obstruction. We first record the fixed-divisor estimate used to detect it.

\section{A prime-power fixed-divisor estimate}


\begin{lemma}
\label{lem:fixed-divisor}

Let $p$ be a prime, let $\nu\ge1$, and put

\nopagebreak[4]
\[
M=p^\nu,
\qquad
V=v_p((M-1)!),
\qquad
P_c(t)=\prod_{r=1}^{M-1}(t-(c+r)).
\]

Then

\nopagebreak[4]
\[
v_p(P_c(t))=
\begin{cases}
V,&t\equiv c\pmod M,\\
\ge V+1,&t\not\equiv c\pmod M.
\end{cases}
\]

\begin{proof}

The assertion is translation-invariant, so write $d=t-c$. If $M\mid d$,
then $v_p(d-r)=v_p(r)$ for $1\le r<M$, giving $V$.

Otherwise choose $h\in\{1,\ldots,M-1\}$ with $d\equiv h\pmod M$. The
factor $d-h$ contributes at least $\nu$, while comparison of the two
remaining factorial segments gives

\nopagebreak[4]
\[
v_p(P_c(t))
\ge \nu+v_p((h-1)!)+v_p((M-1-h)!).
\]

Because $\binom{M-1}{h}$ is nonzero modulo $p$,

\nopagebreak[4]
\[
v_p((h-1)!)+v_p((M-1-h)!)=V-v_p(h).
\]

Since $v_p(h)\le\nu-1$, the result follows.
\end{proof}

\end{lemma}


Set $K=V+1$. Notice that $K\ge\nu$: this is immediate for $\nu=1$, and
for $\nu\ge2$ one has

\nopagebreak[4]
\[
V\ge\left\lfloor\frac{p^\nu-1}{p}\right\rfloor
=p^{\nu-1}-1\ge\nu-1.
\]

\begin{lemma}[Split antipodal obstruction]
\label{lem:split-obstruction}

Let $p\equiv1\pmod4$, let $\nu=v_p(m)\ge1$, and choose
$s\in\mathbb Z/p^\nu\mathbb Z$ with $s^2=-1$. If for some $c$ and $u\in U$

\nopagebreak[4]
\[
A_p(c-s)=\{u\},
\qquad
A_p(c+s)=\{-u\},
\]

then $S_w$ is not a ringset.

\begin{proof}

Let $M=p^\nu$, let $V,K$ be as above, and Hensel-lift $s$ to a square root
of $-1$ modulo $p^K$. In

\nopagebreak[4]
\[
\mathbf L/p^K\mathbf L\simeq M_2(\mathbb Z/p^K\mathbb Z),
\]

put $E=e_u^+$ and $A=1-E=e_u^-$. Choose a matrix $B$ such that

\nopagebreak[4]
\[
ABE\ne0\pmod p.
\]

For example, in a basis where $E=E_{11}$ one may take
$A=E_{22}$ and $B=E_{21}$.

Consider the matrix-coefficient polynomial

\nopagebreak[4]
\[
F(t)=P_c(t)A.
\]

At a point $a+v$, its idempotent decomposition is

\nopagebreak[4]
\[
F(a+v)
=P_c(a+s)Ae_v^+ +P_c(a-s)Ae_v^-.
\]

By Lemma 5.1, a term can survive modulo $p^K$ only if its scalar parameter
is congruent to $c$ modulo $M$. The plus case then has
$a\equiv c-s$ and $v=u$, so $Ae_u^+=AE=0$. The minus case has
$a\equiv c+s$ and $v=-u$, so
$e_{-u}^-=e_u^+=E$ and is also killed. Thus $F$ is null on the periodic
image modulo $p^K$.

After coefficientwise right multiplication by $B$, evaluation at the
appearing point $c-s+u$ has plus component

\nopagebreak[4]
\[
P_c(c)ABE.
\]

Its $p$-adic valuation is exactly $V=K-1$, because $ABE$ is nonzero modulo
$p$. Hence the right multiple is not null modulo $p^K$. Lemma 2.1 proves
that $S_w$ is not a ringset.
\end{proof}

\end{lemma}


\section{Nonsplit prime-power obstructions}


Let $p\equiv3\pmod4$, and suppose that
$A_p(c)=\{q\}$ for some $q\in U$. Put $M=p^\nu$, where
$\nu=v_p(m)$, and retain $V,K$ from Section 5. Define central polynomials

\nopagebreak[4]
\[
\mu_{c+r}(x)=(x-(c+r))^2+1,
\qquad
Q_c(x)=\prod_{r=1}^{M-1}\mu_{c+r}(x),
\]

and

\nopagebreak[4]
\[
h_c(x)=(x-(c+q))Q_c(x).
\]

\begin{lemma}[Higher nonsplit obstruction]
\label{lem:nonsplit-obstruction}

The polynomial $h_c$ is null on the periodic image modulo $p^K$, but a
suitable coefficientwise right multiple is not null. Consequently $S_w$ is
not a ringset.

\begin{proof}

At $y=a+u$, with $d=a-c$, one has

\nopagebreak[4]
\[
\mu_{c+r}(y)=(d-r)(d-r+2u).
\]

The second factor is a unit modulo $p$, because its norm is
$(d-r)^2+4$, which cannot vanish when $-1$ is nonsquare modulo $p$. Hence

For a Lipschitz quaternion $z$, write $v_p(z)$ for the minimum of the
$p$-adic valuations of its four integer coordinates, equivalently for the
largest $s$ such that $z\in p^s\mathbf L$. Then

\nopagebreak[4]
\[
v_p(Q_c(y))=v_p(P_c(a)).
\]

If $M\nmid d$, Lemma 5.1 makes $Q_c(y)$ zero modulo $p^K$. If $M\mid d$,
the periodic orbit lemma says that the only occurring letter is $q$; the
leading factor then evaluates to the real integer $d$, whose valuation is
at least $\nu$. Since $V+\nu\ge K$, $h_c(y)=0$ in every case.

Choose a basis unit $v$ that does not commute with $q$. Since $Q_c$ is
central,

\nopagebreak[4]
\[
h_cv=((x-(c+q))v)Q_c.
\]

At the appearing point $c+q$, the first factor evaluates to

\nopagebreak[4]
\[
v(c+q)-(c+q)v=vq-qv,
\]

which is twice a quaternion unit and hence a $p$-adic unit. The second
factor has valuation exactly $V=K-1$. Thus $h_cv$ is not null modulo $p^K$.
Apply Lemma 2.1.
\end{proof}

\end{lemma}


A constant word may be displayed with period $3$. Taking $p=3$ and $\nu=1$
gives $A_3(c)=\{q\}$ for every position class $c$, so Lemma 6.1 supplies a
constant-word obstruction modulo $3$. Therefore nonconstancy in Theorem 1.1
is necessary even when the originally displayed period has no prime divisor
requiring a separate condition.

\section{The ramified prime \texorpdfstring{$2$}{2}}


The proof at $2$ differs because $-1$ has ramified rather than split or
nonsplit behaviour. Let $\nu=v_2(m)\ge1$, and suppose that one position class
modulo $2^\nu$ contains only a letter $q\in U$. Choose

\nopagebreak[4]
\[
N\ge\max\{\nu,3\},
\qquad
M=2^N,
\]

and refine the monochromatic class to a class $c\pmod M$. Every occurring
point whose real part is congruent to $c$ modulo $M$ still has letter $q$.
Put

\nopagebreak
\[
\begin{gathered}
Q_c(x)=\prod_{r=1}^{M-1}\bigl((x-(c+r))^2+1\bigr),\\[0.4em]
h_c(x)=(x-(c+q))Q_c(x)^2.
\end{gathered}
\]

The square on the central factor is the essential ramified correction.

For $u\in U$, work in the Gaussian subring $\mathbb Z[u]$ and write
$\pi=1+u$. Since $2$ is associated with $\pi^2$, divisibility of all four
Lipschitz coordinates by $2^k$ is equivalent, for elements of
$\mathbb Z[u]$, to $\pi$-adic valuation at least $2k$.

Let

\nopagebreak[4]
\[
P(d)=\prod_{r=1}^{M-1}(d-r),
\qquad
V=v_2((M-1)!).
\]

At $y=a+u$, with $d=a-c$, direct calculation gives

\nopagebreak[4]
\[
Q_c(y)=P(d)P(d+2u).
\]

Define

\nopagebreak[4]
\[
\tau(n)=v_2(n^2+4).
\]

Then

\nopagebreak[4]
\[
v_\pi(Q_c(y))
=2v_2(P(d))+\sum_{r=1}^{M-1}\tau(d-r).
\]

The value of $\tau$ depends only on $n\pmod4$:

\nopagebreak[4]
\[
\tau(n)=
\begin{cases}
0,&n\text{ odd},\\
3,&n\equiv2\pmod4,\\
2,&n\equiv0\pmod4.
\end{cases}
\]

If $M\mid d$, the baseline valuation is

\nopagebreak[4]
\[
L_0=2V+\frac{5M}{4}-2.
\]

Indeed, among $1,\ldots,M-1$ there are $M/2$ odd residues,
$M/4$ residues congruent to $2$ modulo $4$, and $M/4-1$ nonzero residues
divisible by $4$.

\begin{lemma}[Ramified valuation gap]
\label{lem:ramified-gap}

If $M\nmid d$, then

\nopagebreak[4]
\[
v_\pi(Q_c(a+u))\ge L_0+2.
\]

\begin{proof}

Choose $s\in\{1,\ldots,M-1\}$ with $d\equiv s\pmod M$. The consecutive
integer product satisfies

\nopagebreak[4]
\[
v_2(P(d))\ge V-v_2(s)+N.
\]

Compared with the baseline residue set, the $\tau$-sum gains $2$ if $s$ is
odd, loses $1$ if $s\equiv2\pmod4$, and is unchanged if
$s\equiv0\pmod4$. Consequently the total gain is at least

\nopagebreak[4]
\[
2N+2,
\qquad
2N-3,
\qquad
2,
\]

respectively. The last estimate uses $v_2(s)\le N-1$, and the second is at
least $3$ because $N\ge3$.
\end{proof}

\end{lemma}


\begin{lemma}[Ramified monochromatic obstruction]
\label{lem:ramified-obstruction}

If a highest $2$-power position class is monochromatic, then $S_w$ is not a
ringset.

\begin{proof}

Set $K=L_0+2$. Since $K\ge\nu$, the periodic orbit in the quotient modulo
$2^K$ sees the original class modulo $2^\nu$.

If $M\nmid d$, Lemma 7.1 gives

\nopagebreak[4]
\[
v_\pi(Q_c(y)^2)\ge2(L_0+2),
\]

so every Lipschitz coordinate of $Q_c(y)^2$ is divisible by $2^K$. If
$M\mid d$, the occurring letter is $q$ and the leading factor of $h_c(y)$
is the real integer $d$, divisible by $2^N$. The coordinate valuation of
$Q_c(y)^2$ is $L_0$, so again $h_c(y)=0\pmod {2^K}$. Thus $h_c$ is null on
the entire periodic image.

Choose a basis unit $v$ transverse to $q$. Since $Q_c^2$ is central,

\nopagebreak[4]
\[
h_cv=((x-(c+q))v)Q_c^2.
\]

At $c+q$, the leading factor is the commutator $vq-qv$, whose coordinate
$2$-valuation is exactly $1$. The coordinate valuation of
$Q_c(c+q)^2$ is exactly $L_0$. Therefore

\nopagebreak[4]
\[
v_2((h_cv)(c+q))=L_0+1<K,
\]

and the right multiple is not null modulo $2^K$. Lemma 2.1 completes the
proof.
\end{proof}

\end{lemma}


Together with Lemma 3.1, Lemma 7.2 gives the exact local condition at $2$:
all residue images modulo powers of $2$ are core if and only if every
highest $2$-power position class contains at least two letters.

\section{Proof of the classification}

\begin{proof}[Proof of Theorem~\ref{thm:main}]

We now prove Theorem 1.1.

\medskip
\noindent\textbf{Sufficiency.}

Assume conditions 1--3. By Lemma 2.1 and the Chinese remainder theorem, it
is enough to prove that the periodic image is core modulo every prime power
$p^e$.

Suppose first that $p=2$. If $p\nmid m$, every real residue fibre traverses
the entire word by Lemma 2.3 and therefore contains at least two letters by
nonconstancy. If $p\mid m$, condition 2 makes every highest position class
contain at least two letters, and hence every coarser real fibre does also.
Lemma 3.1 applies.

The same argument applies when $p\equiv3\pmod4$: use nonconstancy when
$p\nmid m$, and condition 2 when $p\mid m$. Again Lemma 3.1 gives core
images.

Now let $p\equiv1\pmod4$. If $p\nmid m$, both fibres used at each scalar
parameter in Lemma 4.1 traverse the full nonconstant word, so each contains
at least two letters. If $p\mid m$ and $e\ge\nu_p$, their letter sets are
the highest position fibres in condition 3, after projecting the chosen
square root modulo $p^e$ to modulo $p^{\nu_p}$. Lemma 4.1 applies directly.

If $e<\nu_p$ and the two coarse fibres were singleton antipodes, every
highest-level class refining the first would carry the same singleton, and
every class refining the second would carry its antipode. Hensel-lift the
root modulo $p^e$ to the unique root modulo $p^{\nu_p}$ and lift the scalar
centre arbitrarily. This would violate condition 3. Hence Lemma 4.1 applies
for every $e$.

For an arbitrary $n=\prod_rp_r^{e_r}$, the Chinese remainder theorem gives

\nopagebreak[4]
\[
\mathbf L/n\mathbf L
\simeq\prod_r\mathbf L/p_r^{e_r}\mathbf L.
\]

Polynomial coefficients and evaluation are componentwise, so the null ideal
of the image is the product of the null ideals of its projections. Every
prime-power projection is core, and Lemma 2.1 proves that $S_w$ is a
ringset.

\medskip
\noindent\textbf{Necessity.}

If $w$ is constant, display it with period $3$ and take $p=3$, $\nu=1$.
Every position fibre is then the singleton $\{q\}$, so Lemma 6.1 shows that
$S_w$ is not a ringset.

If condition 2 fails for a prime $p\equiv3\pmod4$, Lemma 6.1 gives an
explicit null polynomial and a failed right multiple. If it fails for
$p=2$, Lemma 7.2 gives the corresponding ramified witness. In both cases
Lemma 2.1 shows that $S_w$ is not a ringset.

Finally, if condition 3 fails at a split prime, Lemma 5.2 gives an explicit
matrix-coefficient null polynomial whose coefficientwise right multiple is
not null. Thus $S_w$ is again not a ringset. All three conditions are
necessary, completing the proof.
\end{proof}


\section{Consequences and examples}


\begin{corollary}[Binary words]
\label{cor:binary}

Let $w:\mathbb Z\to\{\mathbf i,\mathbf j\}$ have arbitrary period $m$.
Then $S_w$ is a ringset if and only if $w$ is nonconstant and, for every
$p\mid m$ with $p=2$ or $p\equiv3\pmod4$, each position class modulo
$p^{v_p(m)}$ is bichromatic.

\begin{proof}

The binary alphabet has no antipodal pair, so condition 3 of Theorem 1.1 is
automatic.
\end{proof}

\end{corollary}


Corollary 9.1 strictly extends \cite{carptopus-odd}, which assumes that $m$ is odd. In
particular, the condition at $2$ is not a formal copy of the odd-prime
condition: its necessity comes from the ramified Gaussian calculation in
Section 7.

\begin{corollary}[Antipode-free alphabets]
\label{cor:antipode-free}

Let $\Omega\subseteq U$ contain at least two letters and no pair $u,-u$.
For a periodic word $w:\mathbb Z\to\Omega$, the set $S_w$ is a ringset if
and only if conditions 1 and 2 of Theorem 1.1 hold.

\end{corollary}

The split condition disappears because no forbidden antipodal singleton pair
can occur. This includes the natural three-letter alphabet
$\{\mathbf i,\mathbf j,\mathbf k\}$.

The full alphabet exhibits genuinely new split failures. If
$p\equiv1\pmod4$ divides the period and two square-root-separated position
classes are constantly $u$ and $-u$, respectively, then all ramified and
nonsplit colour conditions may hold while the set still fails to be a
ringset. Lemma 5.2 detects exactly this configuration.

\section{Computational calibration}


The proof is symbolic and does not depend on finite enumeration. Computation
was used during discovery to locate and attack the local boundaries.

\begin{itemize}
\item All $81$ three-state fibre patterns modulo $4$ and all $256$ binary words
of period $8$ were checked against the evaluation kernel in degrees
strictly below $2q$ for the tested modulus $q$.
\item The $16$ period-$8$ patterns surviving modulo $16$ all failed modulo $64$,
disproving the provisional half-period-antipodal sufficiency guess.
\item The ramified valuation formula was independently checked for
$N=3,\ldots,10$; this exposed that the smallest off-class gain is $2$, not
$3$, and the proof above uses the corrected value.
\item The split eigenline table was checked for all $36$ ordered letter pairs;
exactly the six antipodal pairs fail complementarity.
\item The split obstruction was evaluated for $p=5,13,17$ and
$\nu=1,2$ in finite quotients.

\end{itemize}

These calculations are calibration and destructive testing only. They do
not establish the all-period theorem, whose general content is proved in
Sections 3--8. The first two bullets share the same truncated-null-module
solver and therefore are not independent implementations.

\section{Prior work and novelty boundary}


Frisch \cite{frisch2018} established a general connection between ringsets and null-ideal
sets. Werner \cite{werner2022} classified null ideals of subsets of matrix rings over
fields, while Rissner and Werner \cite{rissnerwerner2025} studied their enumeration. Werner \cite{werner2026}
classified finite Lipschitz and Hurwitz quaternion ringsets and proved the
residue and right-multiplication criteria used here. Peruginelli \cite{peruginelli2014}
developed prime-power fixed-divisor ideals in the commutative setting. None
of these general tools is claimed as new.

The preceding squarefree-period paper \cite{carptopus-squarefree} supplies the initial periodic
framework, and \cite{carptopus-odd} proves the complete all-odd-period binary classification.
The present theorem strictly contains both published results. Its candidate
new content is limited to:

\begin{enumerate}
\item the explicit ramified $2$-adic monochromatic-fibre obstruction;
\item the exact split-prime antipodal singleton obstruction;
\item their assembly into the complete six-letter, arbitrary-period
classification.

\end{enumerate}

The paired-fibre calculation, split matrix model, idempotent vocabulary,
consecutive-product fixed-divisor estimates, and Chinese-remainder assembly
are treated as prior tools or standard ingredients.

As of the bounded search completed on 24 August 2026, no public source was
found that directly states or implies the complete criterion of Theorem 1.1.
This is a qualified \texttt{NO\_DIRECT\_COVERAGE\_FOUND} statement, not a certificate
of absolute global priority. The manuscript has not received external
expert review or formal peer review.

\section*{Declaration of generative AI and AI-assisted technologies}


During the research and preparation of this work, the author used OpenAI
Codex to assist with literature searches, symbolic exploration, exploratory
code review, proof checking, and language drafting. The author reviewed and
edited all generated output and takes full responsibility for the content of
this manuscript.

\begin{thebibliography}{9}
\footnotesize

\bibitem{carptopus-squarefree}
Carptopus,
\emph{Periodic infinite ringsets in the Lipschitz quaternions},
Public Beta v0.1-beta (2026).
\href{https://doi.org/10.5281/zenodo.22069614}{doi:10.5281/zenodo.22069614}.

\bibitem{carptopus-odd}
Carptopus,
\emph{Periodic two-letter ringsets in the Lipschitz quaternions: the general odd-period classification},
Public Beta v0.1-beta (2026).
\href{https://doi.org/10.5281/zenodo.22070532}{doi:10.5281/zenodo.22070532}.

\bibitem{peruginelli2014}
G.~Peruginelli,
\emph{Primary decomposition of the ideal of polynomials whose fixed divisor is divisible by a prime power},
J. Algebra 398 (2014), 227--242.
\href{https://doi.org/10.1016/j.jalgebra.2013.09.016}{doi:10.1016/j.jalgebra.2013.09.016};
\href{https://arxiv.org/abs/1304.7450}{arXiv:1304.7450}.

\bibitem{frisch2018}
S.~Frisch,
\emph{Polynomial functions on non-commutative rings---a link between ringsets and null-ideal sets},
ITM Web Conf. 20 (2018), 01003.
\href{https://doi.org/10.1051/itmconf/20182001003}{doi:10.1051/itmconf/20182001003};
\href{https://arxiv.org/abs/1809.09160}{arXiv:1809.09160}.

\bibitem{werner2022}
N.~J. Werner,
\emph{Null ideals of subsets of matrix rings over fields},
Linear Algebra Appl. 642 (2022), 50--72.
\href{https://doi.org/10.1016/j.laa.2022.02.007}{doi:10.1016/j.laa.2022.02.007}.

\bibitem{werner2026}
N.~J. Werner,
\emph{Integer-valued polynomials on subsets of quaternion algebras},
J. Algebra 686 (2026), 195--219.
\href{https://doi.org/10.1016/j.jalgebra.2025.08.011}{doi:10.1016/j.jalgebra.2025.08.011};
\href{https://arxiv.org/abs/2412.20609}{arXiv:2412.20609v2}.

\bibitem{rissnerwerner2025}
R.~Rissner and N.~J. Werner,
\emph{Counting core sets in matrix rings over finite fields},
Linear Algebra Appl. 720 (2025), 1--25.
\href{https://doi.org/10.1016/j.laa.2025.04.006}{doi:10.1016/j.laa.2025.04.006};
\href{https://arxiv.org/abs/2405.04106}{arXiv:2405.04106v2}.

\end{thebibliography}

\end{document}
